\section{Thảo luận}\label{sec:discussion}

\subsection{Tiến bộ trong mô phỏng thực chất đo cái gì}

Phát hiện trung tâm là: trên dữ liệu lập trình thật, tiến bộ của một chính sách bên trong bộ mô phỏng KT chủ yếu đo mức học mà mô hình của bộ mô phỏng \emph{tin rằng} các lựa chọn bài của nó tạo ra, chứ không phải mức người học sẽ học được. \Cref{prop:insim} phát biểu điều này chính xác: khi một mô hình KT vừa là bộ mô phỏng vừa là thước đo, tiến bộ báo cáo kỳ vọng của mọi chính sách là độ trôi tích lũy của nó, và phần bất ngờ đo được trong DKT-A quả thực không đáng kể. Sáu mô hình người học có độ chính xác dự báo tương đương lại hàm ý mức học rất khác nhau: các mô hình có học tường minh gán cho việc luyện tập những tiến bộ nhỏ và đồng đều, trong khi hai DKT gán tiến bộ vài điểm phần trăm cho từng bài riêng lẻ, theo những hướng mà một seed ngẫu nhiên thứ hai chỉ tái lập yếu. Đây là một trường hợp của tính thiếu xác định (underspecification) \citep{damour2022underspecification}: dữ liệu ràng buộc dự báo bước kế tiếp — điều mọi mô hình đều khớp — nhưng không ràng buộc hệ quả phản thực tế của những chuỗi luyện tập mà người học chưa từng đi qua. Kết quả này mở rộng nhận xét rằng các mô hình người học có độ chính xác tương tự có thể kê đơn lượng luyện tập khác nhau \citep{lee2012individualizing,rollinson2015predictive} từ quyết định về mức thành thạo sang các chính sách sắp xếp trình tự được tối ưu, nơi sự bất đồng bị khuếch đại chứ không chỉ lộ ra. Học tăng cường là một cách tìm kiếm hiệu quả chính những chuỗi như vậy. Bộ tối ưu mạnh hơn (PPO) khai thác phần thưởng DKT nhiều nhất, đạt tiến bộ trong mô phỏng khoảng gấp đôi DQN và một tiến bộ ảo có cùng độ lớn trong thế giới không học.

Hai quan sát nữa làm rõ cách diễn giải. Thứ nhất, với một chính sách cho trước, tiến bộ mà tracker báo cáo thay đổi rất ít giữa chín thế giới — với PPO thưởng DKT trong khoảng {{tg_min_ppo_kt}} đến {{tg_max_ppo_kt}} điểm \%, với chọn ngẫu nhiên từ {{tg_min_random}} đến {{tg_max_random}} điểm \%: giá trị của nó chủ yếu do chính sách chọn bài nào quyết định, chứ không do người học mô phỏng phản ứng ra sao. Thứ hai, dọc theo chuỗi tiếp diễn thật của người học, độ trôi của DKT tự triệt tiêu; dự báo của nó dao động mạnh, như đã được ghi nhận với DKT \citep{yeung2018dktplus}, nhưng không lệch có hệ thống ở đó. Vấn đề phát sinh ngoài phân phối dữ liệu, nơi chính sách cộng dồn độ trôi bước này sang bước khác. Kiểm định dự báo trên chuỗi đã ghi log không thể phát hiện kiểu thất bại này; đó là lý do các mô hình không phân biệt được trong \cref{tab:forecast} vẫn có thể bất đồng về chính sách.

\subsection{Liên hệ với khai thác sai số mô hình và tối ưu quá mức phần thưởng}

Tiến bộ ảo là phiên bản giáo dục của hiện tượng khai thác sai số mô hình trong RL dựa trên mô hình và của việc tối ưu quá mức phần thưởng: một tác tử được tối ưu trên mô hình học từ dữ liệu sẽ tìm ra những hành động mà mô hình lạc quan nhất \citep{kurutach2018metrpo,janner2019trust,gao2023scaling,skalse2022reward}. Bộ gợi ý KT--RL đặc biệt dễ tổn thương vì cùng một mô hình học từ dữ liệu cung cấp động học, phần thưởng và phép đánh giá, nên đánh giá bên trong mô hình đó không thể lộ ra sự khai thác. Các biện pháp phát triển trong RL dựa trên mô hình gợi ý những thiết kế có thể kiểm định cho bộ gợi ý giáo dục: phạt mức bất đồng của một tổ hợp tracker, mà bản song sinh của chúng tôi là trường hợp tối thiểu \citep{kurutach2018metrpo,yu2020mopo}; giới hạn lập kế hoạch ở chân trời ngắn hoặc ở các bài có nhiều dữ liệu \citep{janner2019trust}; hoặc ràng buộc mô hình người học để độ trôi diễn giải được và đơn điệu \citep{zhang2023cmkt}. Chúng tôi chưa đánh giá các biện pháp này, và liệu chúng có khôi phục được tiến bộ học thật hay không vẫn là câu hỏi mở.

\subsection{Hàm ý cho việc đánh giá hệ gợi ý thích ứng}

Kết quả của chúng tôi không cho thấy các hệ thống đã công bố mà chúng tôi dựa vào là sai; chúng tôi không tái lập nhiệm vụ, phần thưởng hay dữ liệu của họ. Kết quả cho thấy một thiết kế chung của nhiều hệ thống — huấn luyện và chấm điểm chính sách bằng cùng một mô hình KT đã fit, hoặc bằng các mô hình fit trên cùng log — có thể báo cáo cải thiện lớn cho những chính sách không dạy được gì. Vì vậy chúng tôi đề xuất bốn kiểm tra ít tốn kém cho các nghiên cứu gợi ý bài tập dựa trên mô phỏng. (i) Báo cáo \emph{đối chứng âm}: đánh giá mọi chính sách trong một thế giới không có học, nơi tiến bộ thật được biết là bằng 0, và báo cáo tiến bộ mà tracker tuyên bố ở đó. (ii) Báo cáo \emph{phép phân rã độ trôi} ở phương trình \eqref{eq:decomposition}, chỉ cần tracker và phơi bày những tiến bộ mà mô hình tự kỳ vọng. (iii) Đánh giá trong các \emph{thế giới fit độc lập} với cấu trúc học và người học khác nhau, theo ma trận đánh giá bền vững \citep{doroudi2017rem}, gồm ít nhất một thế giới có tham số học tường minh. (iv) Fit lại tracker với seed khác và báo cáo mức đồng thuận về mức học ngầm định. Không kiểm tra nào trong số này cần dữ liệu học sinh ngoài log đã dùng. Vì kết quả RL thay đổi theo seed huấn luyện, các kiểm tra nên được báo cáo theo từng seed hoặc kèm độ bất định ở mức seed \citep{henderson2018matters,agarwal2021precipice}.

\subsection{Học tăng cường có thể đóng góp gì}

Học tăng cường không vô dụng. Khi thế giới huấn luyện có cấu trúc học tường minh, diễn giải được (AFM), DQN và PPO học được các chính sách chuyển giao sang các thế giới AFM và PFA fit độc lập và sang một thế giới có quy luật học giả định khác, tiệm cận oracle một bước có đặc quyền trong khi tốn ít lượt nộp hơn chọn ngẫu nhiên. Điều này phù hợp với nhận xét của \citet{doroudi2019reward} rằng chính sách rút ra từ mô hình có ràng buộc lý thuyết đáng tin hơn chính sách rút ra từ mô hình thuần dữ liệu không ràng buộc. Tuy vậy bằng chứng còn hạn chế: tiến bộ khiêm tốn — từ {{afm_adv_min}} đến {{afm_adv_max}} điểm \% xác suất đúng post-test kỳ vọng so với chọn ngẫu nhiên sau mười bài trong AFM-B và PFA-B — vẫn là đại lượng mô phỏng, các thế giới ẩn có chung cấu trúc với thế giới AFM, và chính các tác tử này không vượt chọn ngẫu nhiên trong DKT-B. Các luật nhắm xác suất đúng dự báo vừa phải (0,7) cũng bền vững, trong khi lộ trình dễ$\to$khó và chọn bài dễ nhất trước — thái cực ``quá dễ'' trong lập luận vùng phát triển gần của đề tài — luôn kém hiệu quả nhất trong các thế giới ẩn.

\subsection{Giới hạn}

Mọi kết quả về chính sách đều là mô phỏng. Các thế giới ẩn được fit trên log quan sát của một nền tảng và một nhóm môn học; hệ số học của chúng là các liên hệ đã điều chỉnh theo năng lực và độ khó của bài, không phải hiệu ứng nhân quả, và thế giới ZPD là một giả định. Vì các thế giới ẩn có chung cấu trúc logistic cộng, độ bền của các tác tử huấn luyện với AFM là có điều kiện trên cấu trúc đó. Cách gom kỹ năng là của chúng tôi, và con trỏ — dù nổi bật trong mô tả đề tài — chỉ có một bài trong tập ứng viên vì ít bài con trỏ đạt ngưỡng hỗ trợ 5\%. Công sức được đo bằng số lượt nộp, không phải thời gian làm bài. Chúng tôi nghiên cứu kết quả lần nộp đầu, chân trời mười bài và một trọng số công sức duy nhất; chân trời dài hơn, phần thưởng khác và tập ứng viên khác có thể thay đổi độ lớn kết quả, và chúng tôi chưa phân tích độ nhạy cho các lựa chọn này. Mỗi tác tử được huấn luyện với ba seed, và khoảng tin cậy có điều kiện trên các seed này và trên các mô hình đã fit; chúng không được hiệu chỉnh cho so sánh bội, dù mọi phép so sánh đăng ký trước đều được báo cáo, và các tương quan thứ hạng được tính trên những chính sách không độc lập với nhau. Dữ liệu đến từ hệ thống chấm bài của một trường đại học, nên khả năng khái quát sang ngôn ngữ lập trình, thiết kế môn học và quần thể người học khác chưa được kiểm chứng. Cuối cùng, chúng tôi không thử bất kỳ chính sách nào với học sinh; bằng chứng liên quan tới tính hợp lệ của đánh giá dựa trên mô phỏng và độ bền tương đối của các chính sách trên các thế giới hợp lý, chứ không phải tiến bộ học tập của người học thật.
